\documentclass[11pt]{article}
\usepackage[margin=1in]{geometry}
\usepackage{enumitem}
% =============================================================================
% Preambles/header.tex — shared LaTeX/Beamer preamble for this project
%
% Use from a Beamer lecture by prepending:
%     \input{header}
% and compiling with TEXINPUTS=../Preambles:$TEXINPUTS (the /compile-latex
% skill does this automatically).
%
% PALETTE CONTRACT. Color names defined here MUST match the names in
% Quarto/theme-template.scss so Beamer and Quarto renderings use the same
% palette. The scripts/check-palette-sync.sh script enforces this.
%
% To customize: edit the HEX values below AND the matching $variable in
% Quarto/theme-template.scss, then run scripts/check-palette-sync.sh.
% =============================================================================

% -----------------------------------------------------------------------------
% Engine detection + encoding
%
% Our /compile-latex skill uses XeLaTeX, where inputenc is unnecessary and
% can produce encoding warnings. Only load inputenc under pdfTeX.
% -----------------------------------------------------------------------------
\usepackage{iftex}
\ifPDFTeX
  \usepackage[utf8]{inputenc}
\fi

\usepackage{amsmath, amssymb, amsthm}
\usepackage{booktabs}
\usepackage{graphicx}

% -----------------------------------------------------------------------------
% PALETTE — mirrors Quarto/theme-template.scss variable names.
% Load xcolor and define colors BEFORE anything that references them
% (notably hyperref, hypersetup, and the Beamer color assignments).
% -----------------------------------------------------------------------------
\usepackage{xcolor}

\definecolor{primary-blue}{HTML}{012169}      % headings, accents
\definecolor{primary-gold}{HTML}{B9975B}      % emphasis, borders
\definecolor{highlight-yellow}{HTML}{F2A900}  % markers, alerts
\definecolor{light-bg}{HTML}{E8EDF5}          % title / section backgrounds
\definecolor{jet}{HTML}{1A1A1A}               % body text

% Semantic colors (used in diagrams and annotations)
\definecolor{positive}{HTML}{15803D}          % good / observed / identified
\definecolor{negative}{HTML}{B91C1C}          % bad / problematic / bias
\definecolor{neutral}{HTML}{525252}           % reference / context

% Highlight accents (match .hi-* classes in the SCSS)
\definecolor{hi-slate}{HTML}{314F4F}
\definecolor{hi-green}{HTML}{2E8B57}
\definecolor{hi-red}{HTML}{C41E3A}

% Semantic aliases so skill templates and snippets can write
% \textcolor{accent}{...} without hard-coding "primary-blue".
\colorlet{accent}{primary-blue}
\colorlet{accent2}{primary-gold}

% -----------------------------------------------------------------------------
% Hyperref — loaded AFTER xcolor + palette so link colors resolve.
% -----------------------------------------------------------------------------
\usepackage{hyperref}
\hypersetup{colorlinks=true, linkcolor=primary-blue, urlcolor=primary-blue,
            citecolor=primary-blue, pdfborder={0 0 0}}

% -----------------------------------------------------------------------------
% Course code — set once per deck (after \input{header}, before \begin{document})
% via \coursecode{CS401}. Shown in the footer on every slide so a deck's course
% is identifiable without opening the title page. Empty by default.
% -----------------------------------------------------------------------------
\makeatletter
\newcommand{\coursecode}[1]{\gdef\@coursecodetext{#1}}
\gdef\@coursecodetext{}
\makeatother

% -----------------------------------------------------------------------------
% Beamer theme assignments (only apply under Beamer).
%
% We detect Beamer via \@ifclassloaded{beamer}, which is the documented,
% non-internal way. This must be wrapped in \makeatletter/\makeatother
% because \@ifclassloaded contains an @-catcode character.
% -----------------------------------------------------------------------------
\makeatletter
\@ifclassloaded{beamer}{%
  \usefonttheme{professionalfonts}
  \setbeamercolor{structure}{fg=primary-blue}
  \setbeamercolor{frametitle}{fg=primary-blue}
  \setbeamercolor{title}{fg=primary-blue}
  \setbeamercolor{subtitle}{fg=primary-gold}
  \setbeamercolor{author}{fg=primary-blue}
  \setbeamercolor{institute}{fg=neutral}
  \setbeamercolor{date}{fg=primary-gold}
  \setbeamercolor{itemize item}{fg=primary-blue}
  \setbeamercolor{itemize subitem}{fg=primary-gold}
  \setbeamercolor{alerted text}{fg=primary-gold}
  \setbeamercolor{block title}{fg=white, bg=primary-blue}
  \setbeamercolor{block body}{bg=light-bg}
  % Semantic block variants: block=definition/concept (blue), exampleblock=
  % worked example (green/positive), alertblock=key takeaway or caution (gold).
  % Keep to <=2 colored boxes per slide across ALL three variants (INV-7).
  \setbeamercolor{block title example}{fg=white, bg=positive}
  \setbeamercolor{block body example}{bg=light-bg}
  \setbeamercolor{block title alerted}{fg=white, bg=primary-gold}
  \setbeamercolor{block body alerted}{bg=light-bg}

  % Minimal footer: course code (left) + page X / N (right), no navigation clutter
  \setbeamertemplate{navigation symbols}{}
  \setbeamertemplate{footline}{%
    \color{neutral}\footnotesize\hspace{0.7em}\@coursecodetext\hfill
    \insertframenumber/\inserttotalframenumber\hspace{0.7em}\vspace{0.3em}}
}{}
\makeatother

% -----------------------------------------------------------------------------
% TikZ setup — libraries the snippet gallery uses
% -----------------------------------------------------------------------------
\usepackage{tikz}
\usetikzlibrary{arrows.meta, positioning, calc, decorations.pathreplacing,
                fit, shapes.geometric, backgrounds}

% Shared TikZ styles — snippets and hand-written diagrams can pick these up
% via `\begin{tikzpicture}[every node/.style={...}]` or by naming specific
% styles (e.g., `\node[dag-node] (X) at (0,0) {X};`).
\tikzset{
  % Node styles for causal DAGs and similar
  dag-node/.style = {
    draw=primary-blue, fill=light-bg, rounded corners=2pt,
    minimum width=1.2cm, minimum height=0.9cm, align=center, font=\small
  },
  decision-node/.style = {
    draw=primary-gold, fill=white, diamond, aspect=1.8,
    minimum width=1.8cm, minimum height=1.2cm, align=center, font=\small,
    inner sep=1pt
  },
  % Edge styles — observed vs counterfactual
  observed-edge/.style = {-{Latex[length=2.5mm]}, thick, draw=jet},
  counterfactual-edge/.style = {-{Latex[length=2.5mm]}, thick, draw=neutral, dashed},
  confound-edge/.style = {-{Latex[length=2.5mm]}, thick, draw=negative, dashed},
  % Data-point styles for plots
  observed-dot/.style = {fill=primary-blue, draw=primary-blue, circle, inner sep=1.2pt},
  counterfactual-dot/.style = {fill=white, draw=neutral, circle, inner sep=1.2pt}
}

% -----------------------------------------------------------------------------
% Convenience macros
% -----------------------------------------------------------------------------
\newcommand{\muted}[1]{\textcolor{neutral}{#1}}
\newcommand{\key}[1]{\textcolor{primary-gold}{\textbf{#1}}}
\newcommand{\good}[1]{\textcolor{positive}{#1}}
\newcommand{\bad}[1]{\textcolor{negative}{#1}}

% Standout frame macro: full-bleed dark background for section transitions.
% Beamer-only (uses \begin{frame}/\setbeamercolor/\usebeamerfont) -- do not
% call from an article-class file (e.g. Notes/<CODE>/*-notes.tex). \muted,
% \key, \good, \bad, and \coursecode above are plain xcolor/gdef and are
% safe to use from either class; verified when Notes/article-class support
% was added.
% Expand: \transitionslide{Your transition title}
\newcommand{\transitionslide}[1]{%
  \begingroup
  \setbeamercolor{background canvas}{bg=primary-blue}
  \begin{frame}[plain]
    \centering
    \vfill
    {\usebeamerfont{title}\color{white}\Large #1\par}
    \vfill
  \end{frame}
  \endgroup
}

% Section divider frame (Beamer-only), an alternative to \transitionslide with a
% darker two-line style: near-black (jet) background, a small white label line
% above a larger gold title, and a thin gold rule along the bottom edge.
% Expand: \sectiondivider{Section 2}{Pointers}
\newcommand{\sectiondivider}[2]{%
  \begingroup
  \setbeamercolor{background canvas}{bg=jet}
  \begin{frame}[plain]
    \vspace*{3.0cm}
    {\color{white}\ttfamily\large #1\par}
    \vspace{0.5cm}
    {\color{primary-gold}\LARGE #2\par}
    \begin{tikzpicture}[remember picture, overlay]
      \draw[primary-gold, line width=2.5pt]
        ($(current page.south west)+(0,0.1)$) -- ($(current page.south east)+(0,0.1)$);
    \end{tikzpicture}
  \end{frame}
  \endgroup
}

% -----------------------------------------------------------------------------
% Channel watermark — "Concepts That Click" (YouTube).
%
% Article-class files (CompetitiveExam Q/A sets, Notes, Assignments) get a
% light diagonal draft watermark on every page plus a centered footer naming
% the channel. Beamer decks get a subtle TikZ background watermark instead
% (the footline already carries the course code). Centralized here so every
% MCQ PDF is branded without per-file edits.
% -----------------------------------------------------------------------------
\makeatletter
\@ifclassloaded{article}{%
  \usepackage{draftwatermark}
  \SetWatermarkText{Concepts That Click}
  \SetWatermarkScale{1}
  \SetWatermarkFontSize{2cm}
  \SetWatermarkLightness{0.86}
  \SetWatermarkAngle{45}
  \usepackage{fancyhdr}
  \pagestyle{fancy}
  \fancyhf{}
  \fancyfoot[C]{\footnotesize\textcolor{neutral}{Concepts That Click~|~YouTube: \href{https://www.youtube.com/@concepts.thatclick}{@concepts.thatclick}}}
  \renewcommand{\headrulewidth}{0pt}
  \renewcommand{\footrulewidth}{0pt}
  \fancypagestyle{plain}{%
    \fancyhf{}%
    \fancyfoot[C]{\footnotesize\textcolor{neutral}{Concepts That Click~|~YouTube: \href{https://www.youtube.com/@concepts.thatclick}{@concepts.thatclick}}}%
    \renewcommand{\headrulewidth}{0pt}%
    \renewcommand{\footrulewidth}{0pt}%
  }
}{}
\@ifclassloaded{beamer}{%
  \setbeamertemplate{background}{%
    \begin{tikzpicture}[remember picture, overlay]
      \node[rotate=45, opacity=0.055, align=center]
        at (current page.center)
        {\Huge\textcolor{neutral}{Concepts That Click}};
    \end{tikzpicture}%
  }
}{}
\makeatother 
\coursecode{JKSSB}
\renewcommand{\thesection}{21.\arabic{section}}
\newtheorem{example}{Example}[section]
\newenvironment{definitionbox}[1]{%
  \par\noindent\textbf{Definition (#1).}\ \itshape
}{\par\vspace{0.3em}}
\title{Keyboard Shortcuts --- Detailed Lecture Notes}
\author{JKSSB Computer Awareness}
\date{\today}

\begin{document}
\maketitle

\section{How to read a keyboard shortcut}
A keyboard shortcut is a key, or a combination of keys, that performs a
command directly without opening a menu. Shortcuts save time and reduce the
need for the mouse, and they are a favourite topic in the JKSSB Computer
Awareness paper because the answer is a single exact letter or key. In a
combination written with a plus sign, such as \textbf{Ctrl+C}, the plus means:
hold the first key down, then press the second key. You do not type the plus
sign. Where three keys are written, such as \texttt{Ctrl+Shift+N}, hold the
first two and tap the third.

\begin{definitionbox}{Keyboard shortcut}
A key or key combination that runs a command directly --- the same command a
menu offers --- so the user reaches it from the keyboard instead of with the
mouse.
\end{definitionbox}

The keys used to build combinations are called modifier keys. They are
\textbf{Ctrl} (Control), \textbf{Alt} (Alternate), \textbf{Shift}, and the
\textbf{Windows (Win)} key. A modifier key does very little on its own; it
changes what the next key does. The \textbf{function keys} \textbf{F1} through
\textbf{F12} run along the top row of the keyboard and are usually pressed on
their own, without a modifier.

\begin{center}
\begin{tabular}{p{0.16\textwidth}p{0.28\textwidth}p{0.48\textwidth}}
\toprule
\textbf{Key} & \textbf{Full form} & \textbf{Main role} \\
\midrule
Ctrl & Control & Acts on content: copy, cut, paste, save, find. \\
Alt & Alternate & Acts on windows and menus: switch, close. \\
Shift & Shift & Changes what is typed or selected; builds ranges. \\
Win & Windows key & Opens or controls system features: desktop, Explorer, lock. \\
Fn & Function (laptop) & Gives a second job to the F-keys on many laptops. \\
\bottomrule
\end{tabular}
\end{center}

A practical split helps memory. The \textbf{Ctrl} combinations mostly act on
\emph{content} --- copying, cutting, saving, finding. The \textbf{Alt}
combinations mostly act on \emph{windows} --- switching and closing. The
\textbf{Windows} key opens or controls system-level features such as the
desktop, File Explorer, and the lock screen. Function keys are single keys
whose job depends on the application that is active.

Most shortcut letters are chosen to be easy to remember. \textbf{C} stands for
Copy, and \textbf{X} looks like a cross over the text, so it stands for Cut.
\textbf{V} sits next to X on the keyboard and stands for Paste. \textbf{Z} is
Undo and \textbf{Y} is Redo, the two editing-history keys. \textbf{A} is All
(Select All), \textbf{S} is Save, \textbf{P} is Print, \textbf{F} is Find, and
\textbf{H} opens the find-and-replace box. Learning the letter is more reliable
than memorising a list, because the letter usually hints at the command.

\section{Copy, cut, and paste}
The three most common shortcuts form the clipboard group. The
\textbf{clipboard} is the temporary holding area in memory where copied or cut
items wait before being pasted.

\begin{center}
\begin{tabular}{p{0.20\textwidth}p{0.70\textwidth}}
\toprule
\textbf{Shortcut} & \textbf{Action} \\
\midrule
Ctrl+C & Copy the selected item; the original stays where it is. \\
Ctrl+X & Cut the selected item; the original is removed. \\
Ctrl+V & Paste the copied or cut item at the current position. \\
\bottomrule
\end{tabular}
\end{center}

\begin{definitionbox}{Clipboard}
A temporary area in the computer's memory that holds text, images, or files
after they are copied or cut, until they are pasted.
\end{definitionbox}

The distinction to hold is that \textbf{Copy} and \textbf{Cut} both place the
selection on the clipboard, but only \textbf{Cut} removes the original. A
question that offers \texttt{Ctrl+C} and \texttt{Ctrl+X} as two options is
testing exactly this difference: \texttt{Ctrl+C} leaves the original in place,
and \texttt{Ctrl+X} does not. \texttt{Ctrl+V} only reads the clipboard; it
never decides whether the item was copied or cut.

\begin{example}
A candidate reads, ``Which shortcut copies selected text without removing it
from its original place?'' The answer is \texttt{Ctrl+C}. \texttt{Ctrl+X} also
puts the text on the clipboard, but it deletes the selection from the
document, so it does not satisfy the phrase ``without removing it''.
\end{example}

\section{Undo, redo, and select all}
Three Ctrl shortcuts control the editing history and the current selection.

\begin{itemize}
  \item \textbf{Ctrl+Z} --- Undo the last action. It steps backward through
    recent changes.
  \item \textbf{Ctrl+Y} --- Redo the action that was just undone. It steps
    forward again.
  \item \textbf{Ctrl+A} --- Select All the content of the current window or
    document.
\end{itemize}

Undo and redo are a matched pair and must not be swapped. \texttt{Ctrl+Z}
reverses a change; \texttt{Ctrl+Y} reapplies a change that was reversed.
\texttt{Ctrl+A} does not select one word; it selects everything that can be
selected in the current context. In MS Office the redo key is \texttt{Ctrl+Y};
a few non-Office programs use \texttt{Ctrl+Shift+Z}, but for this exam the
standard answer is \texttt{Ctrl+Y}.

\section{Save, print, find, and replace}
These four shortcuts cover saving a file and working with text inside it.

\begin{center}
\begin{tabular}{p{0.20\textwidth}p{0.70\textwidth}}
\toprule
\textbf{Shortcut} & \textbf{Action} \\
\midrule
Ctrl+S & Save the current file. \\
Ctrl+P & Print. \\
Ctrl+F & Find (open the search box for text). \\
Ctrl+H & Replace (find a word and substitute another). \\
\bottomrule
\end{tabular}
\end{center}

\textbf{Find} and \textbf{Replace} are near neighbours. \texttt{Ctrl+F} only
locates the search text and highlights it; \texttt{Ctrl+H} opens the
find-and-replace function, which can also change every occurrence to a new
word. Similarly, \texttt{Ctrl+S} saves and \texttt{Ctrl+P} prints; the two are
often offered together as distractors. \texttt{Ctrl+S} does not print, and
\texttt{Ctrl+P} does not save.

\begin{example}
A user types a report and wants to change every occurrence of the word
``received'' to ``recd''. \texttt{Ctrl+F} would only find the word and stop
there. The correct shortcut is \texttt{Ctrl+H}, which opens Replace and lets
the user substitute the new text, one occurrence at a time or all at once.
\end{example}

\section{Alt-key combinations}
The Alt key builds the window-management shortcuts.

\begin{itemize}
  \item \textbf{Alt+Tab} --- Switch between open windows or running
    applications. Repeatedly pressing Tab while holding Alt cycles through the
    open windows.
  \item \textbf{Alt+F4} --- Close the active window or application.
\end{itemize}

On the desktop with no application window active, \texttt{Alt+F4} opens the
shutdown or sign-out dialog. The concept to keep clear is that \texttt{Alt+Tab}
\emph{switches} between windows and does not close any of them, while
\texttt{Alt+F4} \emph{closes} the active window. A common wrong answer reverses
the two. Note that \texttt{Alt+F4} closes the whole window; saving the file is
still the job of \texttt{Ctrl+S}.

\section{Function keys}
A function key is a single key, \texttt{F1} to \texttt{F12}, on the top row of
the keyboard. Its action depends on the application that is running. For the
exam, four function keys are high-yield.

\begin{center}
\begin{tabular}{p{0.14\textwidth}p{0.78\textwidth}}
\toprule
\textbf{Key} & \textbf{Standard action} \\
\midrule
F1 & Help. \\
F2 & Rename the selected file or folder (for example in File Explorer). \\
F5 & Refresh the window or folder; in PowerPoint, start the slide show from
     the beginning. \\
F7 & Spelling and grammar check (in MS Word). \\
\bottomrule
\end{tabular}
\end{center}

\begin{definitionbox}{Function key}
One of the keys F1 to F12 on the top row of the keyboard, whose effect is
decided by the active application rather than by a fixed system command.
\end{definitionbox}

In \textbf{PowerPoint}, \texttt{F5} starts a slide show from the beginning and
\texttt{Shift+F5} starts it from the current slide. In \textbf{Windows}, on a
folder or desktop, \texttt{F5} refreshes the view. \texttt{F7} in MS Word
launches the spelling and grammar check. A documented JKSSB item asked the key
that starts a slide show from the beginning, and the answer was \texttt{F5}
(PYQ WO 2026 Q80, provisional key).

\begin{example}
A candidate sees four keys: F1, F2, F5, F7. If the stem asks for Help, the
answer is F1; for Rename, F2; for Refresh or slide show, F5; for the spell
check, F7. Read the stem's verb first, then pick the key that matches that verb.
\end{example}

The table below fixes the near neighbours that JKSSB likes to mix. Each row is
one function key; the middle column is its job; the last column is the trap
(the job a careless candidate might wrongly attach to it).

\begin{center}
\begin{tabular}{p{0.10\textwidth}p{0.32\textwidth}p{0.50\textwidth}}
\toprule
\textbf{Key} & \textbf{Correct job} & \textbf{Common wrong answer} \\
\midrule
F1 & Help & Rename or Refresh \\
F2 & Rename a selected file or folder & Help or New \\
F5 & Refresh; start slide show & Spell check or Save \\
F7 & Spelling and grammar check & Slide show or Help \\
\bottomrule
\end{tabular}
\end{center}

\section{Windows-key shortcuts}
The Windows key (the key with the Windows logo) builds system-level shortcuts.

\begin{center}
\begin{tabular}{p{0.20\textwidth}p{0.70\textwidth}}
\toprule
\textbf{Shortcut} & \textbf{Action} \\
\midrule
Win+D & Show the desktop (minimise all open windows). \\
Win+E & Open File Explorer. \\
Win+L & Lock the computer. \\
\bottomrule
\end{tabular}
\end{center}

Mnemonic: \textbf{D} for Desktop, \textbf{E} for Explorer, \textbf{L} for
Lock. These three are system shortcuts, not content shortcuts, so they do not
depend on which file is open. \texttt{Win+D} does not delete files, and
\texttt{Win+L} does not shut the computer down; it only locks the session, and
the user must sign in again.

The \textbf{Shift} key is a modifier too. It produces capital letters and the
upper symbol of a key, and \texttt{Shift+Click} selects a range of items
between two clicks. \texttt{Shift+F5} starts a PowerPoint slide show from the
current slide. On its own with a letter, Shift does not save, close, or print;
it changes what is typed or selected.

\section{Other high-yield keys}
The shortcuts above are the core of the lesson. A few adjacent keys and
combinations appear often enough in Office documents that they are worth
learning alongside them.

\begin{center}
\begin{tabular}{p{0.22\textwidth}p{0.64\textwidth}}
\toprule
\textbf{Shortcut / key} & \textbf{Action} \\
\midrule
Ctrl+O & Open an existing file. \\
Ctrl+N & Create a new file (new document or window). \\
Ctrl+W & Close the current document or window. \\
Ctrl+B / Ctrl+I / Ctrl+U & Bold / Italic / Underline the selection. \\
Esc & Cancel the current action or close a dialog. \\
Delete & Delete the selected item; in text, erase forward. \\
Backspace & Erase the character to the left of the cursor. \\
Home / End & Move to the start / end of the line. \\
Ctrl+Home / Ctrl+End & Move to the top / bottom of the document. \\
Page Up / Page Down & Move up / down one screen. \\
\bottomrule
\end{tabular}
\end{center}

Two cautions. First, \texttt{Ctrl+O} opens and \texttt{Ctrl+N} creates; do not
confuse the two. Second, \texttt{Ctrl+W} closes a \emph{document}, while
\texttt{Alt+F4} closes the \emph{window or application}; in a browser,
\texttt{Ctrl+W} closes the current tab.

\section{Special keys that stand alone}
A few keys send a command on their own, so they are never written with a plus
sign. They are not combinations, but they are learned together with shortcuts.

\begin{center}
\begin{tabular}{p{0.20\textwidth}p{0.70\textwidth}}
\toprule
\textbf{Key} & \textbf{Action} \\
\midrule
Enter & Confirm a command or start a new line. \\
Esc & Cancel the current action or close a dialog. \\
Tab & Move to the next field, cell, or control. \\
Caps Lock & Toggle capital letters on or off. \\
Backspace & Erase the character to the left of the cursor. \\
Delete & Erase the selected item, or the character to the right. \\
Insert & Toggle Insert and Overtype modes in text editors. \\
Home / End & Move to the start or the end of the current line. \\
\bottomrule
\end{tabular}
\end{center}

\textbf{Caps Lock} and \textbf{Shift} both produce capitals, and this is a
common confusion. Caps Lock changes every letter until it is switched off
again, while Shift changes only the key pressed at that moment. \textbf{Delete}
removes the selection or the character ahead of the cursor; \textbf{Backspace}
removes what is behind the cursor.

\section{Shortcuts in the office desktop workflow}
A short walk-through on the office desktop ties the keys together. A clerk
opens an existing letter with \texttt{Ctrl+O}, types the body, and underlines a
heading with \texttt{Ctrl+U}. The heading is copied with \texttt{Ctrl+C} and
pasted elsewhere with \texttt{Ctrl+V}; a wrong sentence is removed with
\texttt{Ctrl+X}. A mistake is reversed with \texttt{Ctrl+Z}. The clerk searches
the letter for a name with \texttt{Ctrl+F} and changes an old address to a new
one with \texttt{Ctrl+H}. The finished letter is saved with \texttt{Ctrl+S} and
printed with \texttt{Ctrl+P}. While the letter is open, the clerk switches to
the spreadsheet with \texttt{Alt+Tab}, and when the letter is done, the window
is closed with \texttt{Alt+F4}. Later, \texttt{Win+D} shows the desktop,
\texttt{Win+E} opens File Explorer to find last month's folder, and \texttt{Win+L}
locks the machine before the clerk steps away. This single sequence uses almost
every shortcut in the lesson, which is why a scenario item is easy to build
from it.

\section{Exam retrieval and common confusions}
Seven distinctions prevent most errors on this topic.
\begin{enumerate}
  \item \textbf{Ctrl+C} copies and keeps the original; \textbf{Ctrl+X} cuts and
    removes the original. Do not treat them as the same command.
  \item \textbf{Ctrl+Z} is Undo; \textbf{Ctrl+Y} is Redo. Reversing this pair
    is a common wrong answer.
  \item \textbf{Ctrl+F} is Find; \textbf{Ctrl+H} is Replace. Find only locates;
    Replace also substitutes.
  \item \textbf{Ctrl+S} is Save; \textbf{Ctrl+P} is Print. Save stores the
    file; Print sends it to the printer.
  \item \textbf{Alt+Tab} switches between windows; \textbf{Alt+F4} closes the
    active window.
  \item \textbf{F1} is Help, \textbf{F2} is Rename, \textbf{F5} is Refresh or
    slide show, and \textbf{F7} is the spell check (TRAP-32). Do not move a
    job from one key to another.
  \item \textbf{Win+D} shows the desktop, \textbf{Win+E} opens File Explorer,
    and \textbf{Win+L} locks the computer.
\end{enumerate}

\begin{example}
An item asks which key is used to check spelling and grammar in MS Word. Trace
the function keys: F1 is Help, F2 is Rename, F5 is Refresh or slide show, F7 is
the spelling and grammar check. The answer is F7. If the stem instead asked for
Rename, the answer would be F2, and if it asked for Help, the answer would be
F1.
\end{example}

\section{Quick-reference sheet}
Keep this single table for last-minute revision. It gathers every shortcut in
this lesson.

\begin{center}
\begin{tabular}{p{0.15\textwidth}p{0.28\textwidth}p{0.15\textwidth}p{0.28\textwidth}}
\toprule
\textbf{Shortcut} & \textbf{Action} & \textbf{Shortcut} & \textbf{Action} \\
\midrule
Ctrl+C & Copy & Ctrl+S & Save \\
Ctrl+X & Cut & Ctrl+P & Print \\
Ctrl+V & Paste & Ctrl+F & Find \\
Ctrl+Z & Undo & Ctrl+H & Replace \\
Ctrl+Y & Redo & F1 & Help \\
Ctrl+A & Select All & F2 & Rename \\
Alt+Tab & Switch window & F5 & Refresh / slide show \\
Alt+F4 & Close window & F7 & Spell check \\
Win+D & Show desktop & Win+E & File Explorer \\
Win+L & Lock computer & Shift+F5 & Slide show from current \\
\bottomrule
\end{tabular}
\end{center}

\subsection*{Exercises}
\begin{enumerate}[label=\arabic*.]
  \item State the action of Ctrl+C, Ctrl+X, and Ctrl+V.
  \item Distinguish Ctrl+Z from Ctrl+Y, and Ctrl+F from Ctrl+H.
  \item What do Alt+Tab and Alt+F4 do, and how are they different?
  \item Give the standard action of F1, F2, F5, and F7.
  \item What do Win+D, Win+E, and Win+L do? What does Shift+F5 do in
    PowerPoint?
  \item Why is the plus sign in \texttt{Ctrl+C} not typed by the user?
\end{enumerate}

\subsection*{Solutions}
\begingroup\small
\begin{enumerate}[label=\arabic*.]
  \item Ctrl+C copies the selection (the original stays); Ctrl+X cuts it (the
    original is removed); Ctrl+V pastes the copied or cut item at the current
    position.
  \item Ctrl+Z undoes the last action, while Ctrl+Y redoes the action that was
    undone. Ctrl+F opens Find, which only locates text; Ctrl+H opens Replace,
    which can also substitute new text for the found text.
  \item Alt+Tab switches between open windows or applications without closing
    any of them; Alt+F4 closes the active window or application. One switches,
    the other closes.
  \item F1 is Help; F2 renames the selected file or folder; F5 refreshes the
    window (or, in PowerPoint, starts the slide show from the beginning); F7
    runs the spelling and grammar check in MS Word.
  \item Win+D shows the desktop, Win+E opens File Explorer, and Win+L locks
    the computer. Shift+F5 starts a PowerPoint slide show from the current
    slide.
  \item The plus sign is only a way of writing the combination. It means
    ``hold the first key, then press the second''; it is not a character to be
    typed.
\end{enumerate}
\endgroup
\end{document}
